%% kapa: exam=1 round=25 year=2014 period=1 form=A subject=economics subject_ko=경제원론 seq=38 number=78
%% source: Q-Net 감정평가사 기출문제 (cst00309) artlSeq=5101207 file=extracted/1008/5101207/2014년 제25회 감정평가사 1차시험 기출문제/1교시(A형).hwp
%% conversion: kordoc 4.24.51 -> tools/gen1.py (auto); review=proofread-sonnet; flags=none
\begin{question}{78}
자국통화를 지속적으로 저평가(undervaluation)할 때 나타나는 현상으로 옳은 것은?
\begin{choices}
  \item 자국통화의 공급이 감소된다.
  \item 디플레이션이 발생한다.
  \item 국내 재화와 서비스의 가격이 상승한다.
  \item 자국 이자율이 하락한다.
  \item 외환보유고가 고갈된다.
\end{choices}
\end{question}
\answer{3}
